\section{Uvod}

Konverzacijska umjetna inteligencija (AI) i chatbotovi postali su neizostavan dio suvremenog digitalnog ekosustava, transformirajući način na koji poduzeća komuniciraju sa svojim korisnicima. Chatbotovi su dizajnirani kao komunikacijski kanali koji omogućuju neprekidnu dostupnost 24 sata dnevno, 7 dana u tjednu, služeći u različite svrhe, od odgovaranja na često postavljana pitanja do angažiranja korisnika i prikupljanja dubljih uvida u njihove potrebe. Ovi sustavi nastoje oponašati ljudsku interakciju, omogućujući duže i često nestrukturirane razgovore, što ih razlikuje od tradicionalnih sustava za dijalog orijentiranih na zadatke koji se oslanjaju na unaprijed definirane "okvire" ili "slotove" za prikupljanje informacija. Moderni sustavi često kombiniraju aspekte oba pristupa, spajajući fluidnost otvorenog chata sa strukturiranom, ciljanom prirodom dijaloga orijentiranog na zadatke.

Značajan napredak u računalnoj snazi, dostupnost ogromnih skupova podataka i razvoj sofisticiranih modela strojnog učenja doveli su do značajnog poboljšanja sustava za razumijevanje i generiranje prirodnog jezika. Ti su napreci omogućili stvaranje konverzacijskih AI sustava koji mogu generirati odgovore riječ po riječ ili dohvaćati najprikladnije odgovore iz skupa unaprijed definiranih kandidata.

Sveučilišta diljem svijeta suočavaju se s izazovima u pružanju brze i točne podrške studentima, posebno tijekom procesa upisa i tijekom studija. Studenti često postavljaju pitanja o uvjetima upisa, rasporedima predavanja, ispitnim rokovima, stipendijama i drugim administrativnim procesima. Tradicionalni kanali podrške, poput e-pošte, telefonskih poziva ili osobnih konzultacija, često su preopterećeni, što dovodi do kašnjenja i frustracija kod studenata. Osim toga, statične FAQ stranice na web stranicama fakulteta ne mogu uvijek odgovoriti na raznolike oblike pitanja koje studenti postavljaju, posebno ako su formulirana na različite načine.

Chatbotovi temeljeni na umjetnoj inteligenciji nude rješenje za ove izazove, omogućujući automatizaciju odgovora na česta pitanja i pružanje trenutne podrške. Projekt Sembot razvijen je s ciljem stvaranja jednostavnog, skalabilnog i lako prilagodljivog chatbota koji koristi moderne tehnologije za obradu prirodnog jezika. Koristeći Semantic Router i FastAPI, Sembot može razumjeti semantički slične upite i pružiti točne odgovore, dok jednostavno web sučelje omogućuje studentima intuitivnu interakciju.

Ovaj rad pruža sveobuhvatan pregled projekta, uključujući tehničku arhitekturu, implementaciju semantičkog usmjeravanja, upravljanje podacima, prijedloge za sustav preporuke te potencijalne buduće nadogradnje. Također se razmatraju alternative trenutnom pristupu i potencijalni izazovi u implementaciji.